\SourceLeaf{130}
eine billige zu ihrem Vortheil gereichende Bestimmung betrachten, die ihnen die Alternative erspart, entweder schon während der Zeit, wo sie nichts von der Universität empfangen, zur Witwenkasse beitragen, oder später, wenn sie Besoldung erhalten, noch für die ganze Zeit ihres unbesoldeten Professorstandes doppelt nachzahlen zu müssen.

Meine zweite Bemerkung betrifft den Zinsfuss, in Beziehung auf welchen ich dem, was in dem P. M. des U. R. O. gesagt ist, nicht ganz beitreten kann. Mir erscheint vielmehr die Rechnung, nach welcher der jetzige Zinsfuss zu \(4\tfrac1{10}\) proc. ermittelt ist, zum Theil als illusorisch. Ich erkläre mich durch ein Beispiel. Die Oesterreichischen \(4\tfrac12\) proc. Papiere stehen nach dem heutigen Courszettel auf \(103\tfrac58\). Beim Ankauf von einem Banquier wird man, alles eingerechnet, gewiss \textit{über} 104 wirklich zahlen müssen, ich will aber nur bei 104 stehen bleiben. Man erhält also für sein eingezahltes Geld in der Wirklichkeit nur \(4\tfrac{17}{32}\), oder nicht ganz \(4\tfrac13\) proc. Zinsen. Es dauert also wenigstens 12 Jahre, bis man nur sagen kann, dass man wirklich 4 proc. Zinsen genossen hat. Nun werden aber von diesen Papieren alle Jahre \textit{sehr grosse} Summen ausgeloost und zu pari zurückgezahlt. Geschieht die Ausloosung schon nach 2 Jahren, so hat man in der Wirklichkeit nur zusammen \(4\tfrac23\) proc. oder für ein Jahr \(2\tfrac13\) proc. Zinsen genossen, ungerechnet die Kosten, mit welchen jede Einziehung verbunden ist. Für den Besitzer eines solchen Papiers ist es auch immer ein gefährlicher Umstand, dass er, wenn die ihn treffende Ausloosung nicht zu seiner Kenntniss gelangt, er also das Einziehen zu rechter Zeit versäumt, einen sehr bedeutenden Verlust erleiden kann. Für die Witwenkasse wird wohl der Banquier, von dem die Papiere erkauft sind, immer die nöthige Vigilanz ausüben, weil ihm selbst durch jede vorfallende Versur ein Gewinn zuwächst, aber eigentliche Verantwortlichkeit für jeden durch mögliches Uebersehen entstandenen Verlust wird er doch schwerlich auf sich nehmen. In dieser Rücksicht will ich also nicht unterlassen, hiermit die Anzeige zu machen, dass in der heute vor acht Tagen in Wien geschehenen Verloosung von \textit{anderthalb Millionen Gulden} der in Rede stehenden Papiere auch eine der Obligationen der Witwenkasse getroffen ist, nemlich die pag. IX. der Rechnung unter Nr. 52 aufgeführte Litr. P Nr. 15472. Dass ich im Stande bin, diese Anzeige zu machen, verdanke ich nur dem zufälligen Umstande, dass ich heute, wo eben diese Rechnung in meinen Händen ist, die Notiz von der geschehenen Verloosung in einem Zeitungsblatt fand, und mir daher die Designation der ausgeloosten Nummer notirte, um sie zu Hause mit der Capitalliste der Witwenkasse vergleichen zu können, und mit dieser Anzeige will ich denn diese lange Exposition beschliessen.

\noindent 9. Januar 1845.\hfill\textsc{Gauss.}

\pdfbookmark[1]{II. Untersuchung des gegenwärtigen Zustandes der Professorenwitwenkasse zu Göttingen}{witwen-II}
\begin{center}
{\large [II.]\par}\medskip
\textit{Untersuchung des gegenwärtigen Zustandes der Professorenwitwenkasse zu Göttingen.}\par\medskip
\textit{Vorwort.}\par
\end{center}
\noindent\hspace*{\parindent}In dem von mir in der Witwenkassen-Angelegenheit am 8. Januar d. J. abgegebenen Votum habe ich die Methode nach ihren wesentlichen Elementen angedeutet, welche ich für die allein geeignete halte, um zu einem so gründlichen Urtheile, wie die Natur des Gegenstands verstattet, zu gelangen. Ich habe die dort bezeichneten allerdings sehr langwierigen Rechnungen jetzt beendigt, und ihre Resultate sind in der zweiten Abtheilung dieser Denkschrift enthalten.

Da ich jedoch eine nähere Bekanntschaft mit den Grundsätzen derartiger Rechnungen bei den meisten Mitgliedern des Collegiums, welchem diese Schrift vorgelegt wird, nicht voraussetzen darf, so habe ich geglaubt, dass es denselben lieb sein würde, den Gegenstand auch noch von andern Seiten und aus mehr populären Gesichtspunkten erwogen zu sehen. Ist es auch nicht möglich, auf diese Art eigentlich \textit{präcise}
